%%%%%%%%%%%%%%%%%%%%%%%%%%%%%%%%%%%%%%%%%
% Analisi Prompt Engineering
% Documento Professionale con Stile Colorato
% Versione Semplificata - Solo Pacchetti Standard
%%%%%%%%%%%%%%%%%%%%%%%%%%%%%%%%%%%%%%%%%

\documentclass[10pt,twocolumn]{article}

%----------------------------------------------------------------------------------------
%	PACCHETTI
%----------------------------------------------------------------------------------------

\usepackage[utf8]{inputenc}
\usepackage[italian]{babel}
\usepackage[margin=2cm]{geometry}
\usepackage{graphicx}
\usepackage{booktabs}
\usepackage{xcolor}
\usepackage{titlesec}
\usepackage{fancyhdr}
\usepackage{hyperref}
\usepackage{abstract}
\usepackage{enumitem}
\usepackage{amsmath}

%----------------------------------------------------------------------------------------
%	COLORI
%----------------------------------------------------------------------------------------

\definecolor{primarycolor}{RGB}{0,51,102}      % Blu navy
\definecolor{secondarycolor}{RGB}{0,102,204}   % Blu chiaro
\definecolor{accentcolor}{RGB}{220,50,47}      % Rosso accento
\definecolor{darkgray}{RGB}{64,64,64}          % Grigio scuro

%----------------------------------------------------------------------------------------
%	FORMATTAZIONE TITOLI
%----------------------------------------------------------------------------------------

\titleformat{\section}
  {\color{primarycolor}\Large\bfseries}
  {\thesection}{1em}{}
  
\titleformat{\subsection}
  {\color{secondarycolor}\large\bfseries}
  {\thesubsection}{1em}{}
  
\titleformat{\subsubsection}
  {\color{darkgray}\normalsize\bfseries}
  {\thesubsubsection}{1em}{}

%----------------------------------------------------------------------------------------
%	ABSTRACT COLORATO
%----------------------------------------------------------------------------------------

\renewcommand{\abstractnamefont}{\normalfont\Large\bfseries\color{primarycolor}}
\renewcommand{\abstracttextfont}{\normalfont\small}

%----------------------------------------------------------------------------------------
%	HYPERLINK
%----------------------------------------------------------------------------------------

\hypersetup{
    colorlinks=true,
    linkcolor=secondarycolor,
    filecolor=secondarycolor,
    citecolor=primarycolor,
    urlcolor=secondarycolor,
    pdftitle={Analisi Prompt Engineering},
    pdfauthor={Castrese Basile}
}

%----------------------------------------------------------------------------------------
%	HEADER E FOOTER
%----------------------------------------------------------------------------------------

\pagestyle{fancy}
\fancyhf{}
\fancyhead[L]{\small\color{darkgray}Analisi Tecniche Prompt Engineering}
\fancyhead[R]{\small\color{darkgray}Università di Napoli Federico II - A.A. 2025/2026}
\fancyfoot[C]{\small\color{darkgray}\thepage}
\renewcommand{\headrulewidth}{0.5pt}
\renewcommand{\headrule}{\hbox to\headwidth{\color{secondarycolor}\leaders\hrule height \headrulewidth\hfill}}

%----------------------------------------------------------------------------------------
%	INFORMAZIONI DOCUMENTO
%----------------------------------------------------------------------------------------

\title{
    \vspace{-1cm}
    {\color{primarycolor}\Huge\bfseries Analisi e Classificazione Automatica\\[0.3cm]
    delle Tecniche di Prompt Engineering\\[0.2cm]
    per Unit Testing con LLM}
    \vspace{0.5cm}
}

\author{
    \large\textbf{Castrese Basile}\\[0.2cm]
    \normalsize Dipartimento di Ingegneria Elettrica e delle Tecnologie dell'Informazione\\
    \normalsize Università degli Studi di Napoli Federico II\\
    \normalsize Napoli, Italia\\[0.2cm]
    \normalsize\href{mailto:b.castrese@studenti.unisa.it}{b.castrese@studenti.unisa.it}
}

\date{\normalsize Progetto Software Testing - Prof. Porfirio Tramontana\\[0.1cm]A.A. 2025/2026\\[0.2cm]\today}

%----------------------------------------------------------------------------------------

\begin{document}

\twocolumn[
  \begin{@twocolumnfalse}
    \maketitle
    
    \begin{abstract}
    \noindent L'efficacia dei Large Language Models (LLM) nella generazione di codice dipende criticamente dalla qualità e dalla tecnica dei prompt utilizzati. Questo studio analizza \textbf{225 prompt} sviluppati da 75 studenti per generare unit test su tre applicazioni Java, con l'obiettivo di validare le classificazioni fornite dagli studenti e analizzare la distribuzione delle tecniche utilizzate. È stato sviluppato un \textbf{tool web-based in Python (Streamlit)} che classifica automaticamente i prompt utilizzando oltre 150 pattern regex organizzati gerarchicamente in una tassonomia di \textbf{18 tecniche} (Base, Reasoning Enhancement, Code Generation, Metacognizione). Un campione strategico di \textbf{40 prompt (17.8\%)} è stato selezionato per validazione manuale, bilanciato tra casi dubbi (confidence < 70\%), discordanze studente-algoritmo e concordanze. I risultati forniscono insights sulla diffusione delle tecniche avanzate di prompting e sull'accuratezza delle auto-classificazioni studentesche.
    \end{abstract}
    
    \vspace{0.3cm}
    \noindent\textbf{\color{primarycolor}Keywords:} Prompt Engineering, Large Language Models, Unit Testing, Classificazione Automatica, Pattern Matching, Chain-of-Thought
    
    \vspace{0.5cm}
  \end{@twocolumnfalse}
]

%----------------------------------------------------------------------------------------
%	CONTENUTO
%----------------------------------------------------------------------------------------

\section{Introduzione}

L'utilizzo dei Large Language Models (LLM) per la generazione automatica di codice e test unitari rappresenta una delle applicazioni più promettenti dell'intelligenza artificiale generativa nel campo dello sviluppo software. La qualità degli output prodotti dagli LLM dipende criticamente dalla formulazione dei \textit{prompt}, ovvero dalle istruzioni fornite al modello \cite{wei2022chain,kojima2022large}.

Durante un'esercitazione didattica sull'uso di LLM per Unit Testing, \textbf{75 studenti} hanno sviluppato prompt per generare test unitari per tre applicazioni Java:

\begin{itemize}[noitemsep,leftmargin=*]
    \item \textbf{TennisScoreManager}: Sistema di gestione punteggio tennis con regole complesse
    \item \textbf{SubjectParser}: Parser per codici materie universitarie
    \item \textbf{HSLColor}: Convertitore colori da spazio HSL a RGB
\end{itemize}

Ogni studente ha fornito una propria classificazione della tecnica di prompting utilizzata, spesso basata su categorie comuni quali \textit{Zero-Shot Learning}, \textit{Few-Shot Learning} e \textit{Chain-of-Thought}.

Tuttavia, le classificazioni fornite dagli studenti risultano frequentemente incomplete o imprecise, non considerando tassonomie più articolate che distinguono tra tecniche avanzate come \textit{Tree-of-Thought} \cite{yao2023tree}, \textit{Program-of-Thought} \cite{chen2022program}, \textit{Self-Refine}, e altri approcci sofisticati.

\subsection{Obiettivi dello Studio}

Il presente studio si pone due obiettivi principali:

\begin{enumerate}[noitemsep,leftmargin=*]
    \item \textbf{Validare le classificazioni}: Verificare la correttezza e completezza delle classificazioni fornite dagli studenti, utilizzando una tassonomia completa di 18 tecniche organizzate gerarchicamente
    \item \textbf{Analisi statistica}: Produrre statistiche sulla varietà di prompt proposti e sulla frequenza di utilizzo delle diverse tecniche
\end{enumerate}

Il dataset comprende \textbf{225 prompt totali} (75 studenti × 3 applicazioni) con informazioni sul testo completo, l'LLM utilizzato e la classificazione proposta.

%------------------------------------------------

\section{Metodologia}

\subsection{Tool Sviluppato}

Per automatizzare l'analisi è stata sviluppata una \textbf{web application in Python} utilizzando il framework \textit{Streamlit}. Il tool integra:

\begin{itemize}[noitemsep,leftmargin=*]
    \item \textbf{Pandas}: Data processing e elaborazione Excel
    \item \textbf{Plotly}: Visualizzazioni interattive
    \item \textbf{Regex}: Classificazione automatica (150+ pattern)
    \item \textbf{Interfaccia Web}: Validazione manuale interattiva
\end{itemize}

\subsection{Workflow Operativo}

Il processo si articola in \textbf{quattro fasi}:

\subsubsection{Fase 1: Caricamento Dati}

Il tool carica il file Excel da Google Forms ed estrae: testo prompt, LLM utilizzato con normalizzazione (Gemini/gemini/gemin → Gemini), classificazione studente.

\subsubsection{Fase 2: Classificazione Automatica}

L'algoritmo analizza ogni prompt usando pattern regex gerarchici in 4 categorie (Tabella \ref{tab:taxonomy}). La classificazione segue priorità: \textbf{Metacognizione} $>$ \textbf{Code Generation} $>$ \textbf{Reasoning} $>$ \textbf{Base}.

Per ogni prompt si calcola:
\begin{itemize}[noitemsep,leftmargin=*]
    \item \textbf{Confidence Score} (0-100\%)
    \item \textbf{Is\_Doubtful}: True se confidence < 70\%
    \item \textbf{Needs\_Correction}: True se studente ≠ algoritmo
\end{itemize}

\begin{table}[htbp]
\caption{Tassonomia delle tecniche riconosciute}
\label{tab:taxonomy}
\centering
\small
\begin{tabular}{@{}llc@{}}
\toprule
\textbf{Categoria} & \textbf{Tecnica} & \textbf{Punti} \\
\midrule
\multirow{2}{*}{Base} & Zero-Shot & +1 \\
 & Few-Shot & +1 \\
\midrule
Reasoning & CoT / Auto-CoT / L2M & +2 \\
Enhancement & ToT / GoT & +2 \\
 & Self-Consistency / S2A & +2 \\
\midrule
Code & PoT / CoC / Scratchpad & +2 \\
Generation & Scratchpad CoT & +1 \\
\midrule
Metacogn. & Self-Refine / Fix & +2 \\
\bottomrule
\end{tabular}
\end{table}

\subsubsection{Fase 3: Campione Validazione}

Selezione strategica di \textbf{40 prompt (17.8\%)} distribuiti come in Tabella \ref{tab:sample}.

\begin{table}[htbp]
\caption{Composizione campione validazione}
\label{tab:sample}
\centering
\small
\begin{tabular}{@{}lcc@{}}
\toprule
\textbf{Tipo} & \textbf{N} & \textbf{Motivazione} \\
\midrule
Dubbi & 10 & Confidence < 70\% \\
Con Correz. & 15 & S ≠ A \\
Senza Corr. & 15 & S = A \\
\midrule
\textbf{Totale} & \textbf{40} & Bilanciato \\
\bottomrule
\end{tabular}
\end{table}

\textbf{Razionale}: Campione statisticamente significativo (~18\%) ma gestibile per validazione accurata.

\subsubsection{Fase 4: Validazione Manuale}

Interfaccia web per compilare: \textit{Manual\_Validation}, \textit{Agreement} (Studente/Algoritmo/Altro), \textit{Manual\_Notes}.

%------------------------------------------------

\section{Risultati}

\subsection{Statistiche Descrittive}

\textit{[Da compilare dopo analisi file Prompt.xlsx]}

Parametri chiave preliminari:

\begin{itemize}[noitemsep,leftmargin=*]
    \item Totale prompt: \textbf{225}
    \item Correzioni necessarie: \textbf{[VALORE]} (\%)
    \item Casi dubbi: \textbf{[VALORE]} (\%)
    \item Confidence media: \textbf{[VALORE]}\%
    \item Tecniche rilevate: \textbf{[N]}/18
    \item LLM: Gemini [N]\%, ChatGPT [N]\%
\end{itemize}

\subsection{Distribuzione Tecniche}

\textit{[Da completare con foglio "Technique\_Distribution"]}

I risultati preliminari suggeriscono predominanza di \textbf{tecniche base} (Few-Shot, Zero-Shot) rispetto ad approcci sofisticati (ToT, PoT).

%\begin{figure}[htbp]
%\centering
%\includegraphics[width=\linewidth]{figures/technique_distribution}
%\caption{Distribuzione tecniche di prompting}
%\label{fig:techniques}
%\end{figure}

\subsection{Validazione Manuale}

\textit{[Da completare dopo validazione 40 prompt]}

Dei 40 prompt validati:

\begin{itemize}[noitemsep,leftmargin=*]
    \item Accuratezza Algoritmo: [N]/40 ([X]\%)
    \item Accuratezza Studenti: [N]/40 ([X]\%)
    \item Accordo S-A: [X]\%
\end{itemize}

Metriche (Precisione, Recall, F1) calcolate da colonna Agreement.

%------------------------------------------------

\section{Discussione}

\subsection{Casi Dubbi Significativi}

L'analisi rivela pattern di ambiguità ricorrenti. Esempio: prompt con esempi minimali classificati Zero-Shot dagli studenti ma Few-Shot dall'algoritmo.

\paragraph{Caso Esempio}
Prompt: \textit{"Create method to parse codes. Example: CS101 → Dept: CS, Code: 101"}

\begin{itemize}[noitemsep,leftmargin=*]
    \item Studente: \textbf{Zero-Shot}
    \item Algoritmo: \textbf{Few-Shot} (72\% conf.)
    \item Validazione: \textbf{Few-Shot} ✓
\end{itemize}

Motivazione: Presenza esempio esplicito Input→Output configura apprendimento per analogia.

\subsection{Limitazioni}

\begin{itemize}[noitemsep,leftmargin=*]
    \item \textbf{Regex}: Non cattura sfumature semantiche complesse
    \item \textbf{Multi-tecnica}: Alcune prompt hanno caratteristiche multiple
    \item \textbf{Campione}: 17.8\% validato manualmente
    \item \textbf{Dominio}: Specifico per unit testing
\end{itemize}

\subsection{Implicazioni Pedagogiche}

Gli studenti tendono a \textbf{sottovalutare la complessità} delle tecniche usate. Gap evidenziato:

\begin{enumerate}[noitemsep,leftmargin=*]
    \item Formazione su tassonomie avanzate
    \item Esempi che distinguono tecniche
    \item Esercitazioni guidate di identificazione
\end{enumerate}

%------------------------------------------------

\section{Conclusioni}

Studio sistematico di 225 prompt con tool automatico e validazione campionaria.

\subsection{Risultati Principali}

\textit{[Completare dopo analisi]}

\begin{itemize}[noitemsep,leftmargin=*]
    \item Tecnica più usata: \textbf{[NOME]} ([X]\%)
    \item Accordo S-A: \textbf{[X]\%}
    \item Accuratezza algoritmo: \textbf{[X]\%}
    \item Varietà: \textbf{[N]}/18 tecniche
\end{itemize}

\subsection{Contributi}

Il tool è efficace per:

\begin{enumerate}[noitemsep,leftmargin=*]
    \item Classificazione rapida grandi volumi
    \item Identificazione automatica casi dubbi
    \item Interfaccia intuitiva validazione
    \item Export strutturato Excel 7-fogli
\end{enumerate}

\subsection{Lavori Futuri}

\begin{itemize}[noitemsep,leftmargin=*]
    \item ML-based classification (BERT, Sentence-BERT)
    \item Classificazione multi-label per prompt ibridi
    \item Validazione estesa 30-40\% dataset
    \item Inter-rater reliability multipli annotatori
\end{itemize}

%------------------------------------------------

\section{Materiali e Metodi}

\subsection{Dataset e Preprocessing}

Dataset da Google Forms: 75 studenti × 3 prompt = 225 totali. File Excel processato con \textit{pandas} per estrarre testo, LLM (normalizzato), classificazione studente.

Normalizzazione gestisce: case-sensitivity, versioni (Gemini Pro 3 → Gemini), typos (gemin → Gemini), placeholder ("Which LLM?" → Non Specificato).

\subsection{Algoritmo Classificazione}

Pattern matching gerarchico a 4 livelli:

\begin{enumerate}[noitemsep,leftmargin=*]
    \item \textbf{Metacognizione}: \textit{refine, improve, iterate, step back, fix}
    \item \textbf{Code Gen}: \textit{program-of-thought, chain-of-code, scratchpad}
    \item \textbf{Reasoning}: \textit{tree-of-thought, CoT, step-by-step, self-consistency}
    \item \textbf{Base}: \textit{example, e.g.} (Few-Shot); default (Zero-Shot)
\end{enumerate}

Confidence = f(pattern matched, specificità, keywords). Is\_Doubtful se < 70\% o pattern competono.

\subsection{Implementazione}

Python 3.9+ con:

\begin{itemize}[noitemsep,leftmargin=*]
    \item \texttt{streamlit 1.28+}: Web framework
    \item \texttt{pandas 2.0+}: Data processing
    \item \texttt{plotly 5.17+}: Visualizzazioni
    \item \texttt{openpyxl}: Excel I/O
    \item \texttt{re}: Regex (standard library)
\end{itemize}

Codice disponibile su GitHub (link in preparazione).

%----------------------------------------------------------------------------------------
%	RINGRAZIAMENTI
%----------------------------------------------------------------------------------------

\section*{Ringraziamenti}

Si ringrazia il professore del corso di Software Testing per il dataset e gli studenti partecipanti per aver reso possibile questo studio.

%----------------------------------------------------------------------------------------
%	BIBLIOGRAFIA
%----------------------------------------------------------------------------------------

\begin{thebibliography}{9}

\bibitem{wei2022chain}
Wei J, Wang X, Schuurmans D, et al.
\textit{Chain-of-Thought Prompting Elicits Reasoning in Large Language Models}.
Advances in Neural Information Processing Systems 35:24824-24837, 2022.

\bibitem{kojima2022large}
Kojima T, Gu SS, Reid M, Matsuo Y, Iwasawa Y.
\textit{Large Language Models are Zero-Shot Reasoners}.
Advances in Neural Information Processing Systems 35:22199-22213, 2022.

\bibitem{yao2023tree}
Yao S, Yu D, Zhao J, et al.
\textit{Tree of Thoughts: Deliberate Problem Solving with Large Language Models}.
arXiv preprint arXiv:2305.10601, 2023.

\bibitem{chen2022program}
Chen W, Ma X, Wang X, Cohen WW.
\textit{Program of Thoughts Prompting: Disentangling Computation from Reasoning for Numerical Reasoning Tasks}.
arXiv preprint arXiv:2211.12588, 2022.

\end{thebibliography}

%----------------------------------------------------------------------------------------

\end{document}
